% Vietnamese translation for OI Public Library
% Source-File: IOI中国国家候选队论文/2007/day2/7.胡伯涛《最小割模型在信息学竞赛中的应用》.ppt
\documentclass[11pt]{article}
\usepackage{oipl-vn}

\newcommand{\Oh}{\mathcal{O}}
\newcommand{\INF}{\infty}
\newcommand{\xor}{\mathbin{\mathrm{xor}}}

\title{Bản ghi chú trình chiếu: Mô hình cắt nhỏ nhất trong thi tin học}
\author{Hu Botao}
\date{2007}

\begin{document}
\maketitle

\origtitle{Mô hình cắt nhỏ nhất và ứng dụng trong thi tin học}
\sourcepath{IOI China candidate team papers / 2007 / day 2 / Hu Botao.ppt}

\begin{abstract}
Các slide trình bày mô hình cắt nhỏ nhất như một công cụ mô hình hóa bài
toán chọn tập. Trọng tâm là bốn nhóm ứng dụng: dùng trực tiếp định nghĩa cắt,
đồ thị đóng trọng số lớn nhất, đồ thị con mật độ lớn nhất, và tập phủ đỉnh
trọng số nhỏ nhất cùng tập độc lập trọng số lớn nhất trong đồ thị hai phía.
\end{abstract}

\paragraph{Từ khóa.}
Luồng mạng; luồng cực đại; cắt nhỏ nhất; đồ thị đóng trọng số lớn nhất; đồ
thị con mật độ lớn nhất; tập phủ đỉnh; tập độc lập.

\section{Cắt nhỏ nhất}

Trong mạng luồng \(G=(V,E)\) có nguồn \(s\) và đích \(t\), một cắt
\([S,T]\) là một phân hoạch \(V=S\cup T\), \(S\cap T=\varnothing\), với
\(s\in S\) và \(t\in T\). Sức chứa của cắt là
\[
  c[S,T]=\sum_{u\in S}\sum_{v\in T}c(u,v),
\]
tức chỉ tính các cạnh đi từ phía nguồn sang phía đích. Định lý luồng cực
đại--cắt nhỏ nhất cho phép tìm cắt nhỏ nhất bằng một lần tính luồng cực đại.

Trong mô hình hóa, mỗi đỉnh thường biểu diễn một quyết định nhị phân: đặt nó
về phía \(S\) hay phía \(T\). Cạnh \(s\to v\) biểu diễn mất mát nếu không
chọn \(v\); cạnh \(v\to t\) biểu diễn mất mát nếu chọn \(v\). Cạnh dung lượng
rất lớn \(\INF\) dùng để cấm một cấu hình không hợp lệ.

\section{Ứng dụng trực tiếp}

\subsection{Network Wars}

Bài toán yêu cầu tìm một tập cạnh cắt tách \(s\) khỏi \(t\) sao cho trọng số
trung bình của các cạnh bị cắt là nhỏ nhất:
\[
  \min_C\frac{\sum_{e\in C}w_e}{|C|}.
\]
Đây là một quy hoạch phân thức \(0\)-\(1\). Với giá trị đoán \(\lambda\), đổi
trọng số mỗi cạnh thành
\[
  w'_e=w_e-\lambda.
\]
Khi đó ta cần tìm cắt \(s\)-\(t\) có tổng \(w'_e\) nhỏ nhất. Dấu của giá trị
tối ưu cho biết \(\lambda\) đang lớn hay nhỏ hơn đáp án, nên có thể nhị phân
trên \(\lambda\). Cắt nhỏ nhất không xuất hiện trực tiếp trong đề, nhưng là
bài toán con sau khi biến đổi mục tiêu phân thức.

\subsection{Optimal Marks}

Mỗi đỉnh của đồ thị cần được gán một nhãn nguyên không âm; một số nhãn đã
biết. Trọng số cạnh là \(l_u\xor l_v\), và cần tối thiểu hóa tổng trọng số
cạnh. Vì phép \(\xor\) tách theo bit,
\[
  a\xor b=\sum_i2^i(a_i\xor b_i),
\]
ta giải độc lập từng bit.

Với một bit cố định, nhãn chỉ là \(0\) hoặc \(1\). Cạnh đóng góp \(1\) khi
hai đầu khác phía và \(0\) khi cùng phía, đúng bằng hình dạng của một cắt.
Thêm cạnh \(\INF\) từ nguồn hoặc tới đích để ép các đỉnh đã biết vào phía
đúng; các cạnh gốc có dung lượng \(1\). Cắt nhỏ nhất trên mạng này cho giá
trị tối ưu của bit đang xét, rồi cộng kết quả theo trọng số \(2^i\).

\section{Đồ thị đóng trọng số lớn nhất}

Trong đồ thị có hướng \(G=(V,E)\), một tập đỉnh \(V'\) là \term{đóng} nếu
\[
  u\in V',\ (u,v)\in E \quad\Rightarrow\quad v\in V'.
\]
Mỗi đỉnh có trọng số \(w_v\), có thể dương hoặc âm. Bài toán là tìm tập đóng
có tổng trọng số lớn nhất.

Cách dựng mạng:
\[
\begin{aligned}
V_N &= V\cup\{s,t\},\\
E_N &= E
 \cup\{(s,v)\mid w_v>0\}
 \cup\{(v,t)\mid w_v<0\}.
\end{aligned}
\]
Các cạnh gốc của \(E\) có dung lượng \(\INF\); cạnh \(s\to v\) có dung lượng
\(w_v\) nếu \(w_v>0\); cạnh \(v\to t\) có dung lượng \(-w_v\) nếu \(w_v<0\).

Một cắt hữu hạn không thể cắt cạnh \(\INF\). Vì vậy, nếu \(S-\{s\}=V_1\) thì
mọi cạnh đi ra khỏi \(V_1\) vẫn phải ở trong \(V_1\); nói cách khác \(V_1\)
là tập đóng. Với \(V^+=\{v:w_v>0\}\), dung lượng cắt là
\[
  c[S,T]=\sum_{v\in V^+-V_1}w_v+\sum_{\substack{v\in V_1\\w_v<0}}(-w_v).
\]
Do đó
\[
  w(V_1)=\sum_{v\in V^+}w_v-c[S,T].
\]
Tổng trọng số dương là hằng số, nên cắt nhỏ nhất tương ứng với tập đóng trọng
số lớn nhất.

\subsection{Profit}

Trong bài \textit{Profit}, xây trạm \(i\) tốn \(p_i\); nhóm khách hàng \(j\)
chỉ được phục vụ nếu hai trạm \(a_j,b_j\) đều được xây, và khi đó sinh lợi
\(c_j\). Tạo một đỉnh trọng số dương \(c_j\) cho mỗi nhóm khách hàng, một đỉnh
trọng số âm \(-p_i\) cho mỗi trạm, rồi nối nhóm khách hàng \(j\) tới hai trạm
\(a_j,b_j\). Chọn nhóm khách hàng kéo theo chọn hai trạm, đúng là ràng buộc
đóng.

\section{Đồ thị con mật độ lớn nhất}

Mật độ của đồ thị con \(G'=(V',E')\) là
\[
  D'=\frac{|E'|}{|V'|}.
\]
Với giá trị đoán \(g\), xét bài toán
\[
  h(g)=\max_{V',E'}\{|E'|-g|V'|\}.
\]
Nếu \(h(g)>0\) thì tồn tại đồ thị con có mật độ lớn hơn \(g\); nếu \(h(g)<0\)
thì \(g\) vượt quá mật độ tối ưu. Vì hai mật độ khác nhau trong đồ thị \(n\)
đỉnh cách nhau ít nhất \(1/n^2\), chỉ cần \(\Oh(\log n)\) lần nhị phân trong
trường hợp không trọng số.

\subsection{Cách dựng sơ bộ}

Biến mỗi cạnh \(e=(u,v)\) thành một đỉnh \(v_e\) trọng số \(1\), giữ mỗi đỉnh
gốc \(v\) với trọng số \(-g\), rồi thêm hai cung \(v_e\to u\) và \(v_e\to v\).
Chọn một cạnh bắt buộc chọn hai đầu mút của nó, nên bài toán trở thành đồ thị
đóng trọng số lớn nhất. Cách này đúng nhưng tạo mạng kích thước \(n+m\).

\subsection{Cách dựng cải tiến}

Với \(V'\) cố định, luôn nên chọn toàn bộ cạnh trong đồ thị con cảm sinh bởi
\(V'\). Gọi \(d_v\) là bậc của \(v\). Ta có
\[
  |E'|=\frac{\sum_{v\in V'}d_v-c[V',\overline{V'}]}2.
\]
Vì vậy tối đa hóa \(|E'|-g|V'|\) tương đương tối thiểu hóa
\[
  \sum_{v\in V'}(2g-d_v)+c[V',\overline{V'}],
\]
sai khác một hằng số và hệ số cố định.

Dựng mạng trên các đỉnh gốc:
\[
\begin{aligned}
V_N&=V\cup\{s,t\},\\
E_N&=\{(u,v),(v,u)\mid (u,v)\in E\}
\cup\{(s,v),(v,t)\mid v\in V\}.
\end{aligned}
\]
Đặt \(c(u,v)=1\) cho các cung ứng với cạnh gốc, \(c(s,v)=U\), và
\[
  c(v,t)=U+2g-d_v,
\]
với \(U=m\) để mọi dung lượng không âm. Nếu \(V'=S-\{s\}\), dung lượng cắt là
\[
  c[S,T]=Un-2(|E'|-g|V'|).
\]
Do \(Un\) là hằng số, cắt nhỏ nhất cho \(h(g)\). Mạng chỉ còn \(n+2\) đỉnh và
\(\Theta(n+m)\) cạnh.

\section{Tập phủ và tập độc lập trong đồ thị hai phía}

Trong đồ thị hai phía \(G=(X,Y,E)\), tập phủ đỉnh trọng số nhỏ nhất có thể
được giải bằng cắt nhỏ nhất. Dựng mạng:
\[
  s\to x\quad(x\in X),\qquad x\to y\quad((x,y)\in E),\qquad y\to t\quad(y\in Y).
\]
Dung lượng \(s\to x\) là trọng số của \(x\), dung lượng \(y\to t\) là trọng
số của \(y\), còn cạnh \(x\to y\) có dung lượng \(\INF\). Mỗi cạnh gốc
\((x,y)\) phải được phủ bởi ít nhất một đầu mút; cạnh \(\INF\) đảm bảo không
thể có \(x\in S\) và \(y\in T\) đồng thời không trả chi phí phủ.

Phần bù của một tập phủ là một tập độc lập. Vì vậy, nếu tổng trọng số toàn bộ
đỉnh là \(W\), thì
\[
  \text{trọng số tập độc lập lớn nhất}
  =W-\text{trọng số tập phủ nhỏ nhất}.
\]
Các ứng dụng trong slide gồm \textit{Destroying the Graph} và \textit{Exca},
đều dựa vào bước biến ràng buộc bài toán thành phủ đỉnh trên đồ thị hai phía.

\section{Khuôn mẫu và kỹ thuật}

Cách nghĩ chung của mô hình cắt nhỏ nhất là biến lời giải thành lựa chọn hai
phía của từng đỉnh. Lợi ích dương thường nối từ nguồn, chi phí thường nối tới
đích, còn quan hệ ``chọn \(u\) thì phải chọn \(v\)'' dùng cạnh \(\INF\) từ
\(u\) tới \(v\). Khi bài toán có tỉ số, nhị phân tham số rồi chuyển về bài
toán tổng; khi bài toán có \(\xor\), tách theo bit; khi bài toán có ràng buộc
phủ trên đồ thị hai phía, dùng cạnh \(\INF\) để cấm cạnh không được phủ.

\end{document}
